% GENERATED FILE. Edit canonical lessons, then run npm run generate:books.
% canonical-source-hash: fnv1a64:27a87f4e0f8e3aec
% canonical-lessons: KA-C280-pisuguttu, KA-C280-cappale-tattu, KA-C280-kai-kuluku, KA-C280-talupisu, KA-C280-hindikku, KA-R280-first-pass-words-from-chapters-272-275, KA-R280-second-pass-words-from-chapters-276-280

\chapter{To whisper, To clap, To shake hands, To deliver, To overtake}
\label{ch:ka-to-whisper-to-clap-to-shake-hands-to-deliver-to-overtake}
\hlchaptermodality{280}

\begin{chapteropening}
\textbf{By the end of this chapter:} \emph{I can say to whisper, to clap, to shake hands, to deliver and to overtake.}

Complete the last lesson of chapter 280: I can say to whisper, to clap, to shake hands, to deliver and to overtake.
\end{chapteropening}

\section[\texorpdfstring{\kn{ಪಿಸುಗುಟ್ಟು} (pisuguṭṭu)}{ಪಿಸುಗುಟ್ಟು (pisuguṭṭu)}]{\kn{ಪಿಸುಗುಟ್ಟು} (pisuguṭṭu) — to whisper}

\label{lesson:KA-C280-pisuguttu}



\begin{quote}
Before the new one: say the Kannada for to solve, then the Kannada for to dig.
\end{quote}

\subsection*{You'll want to know: \kn{ಪಿಸುಗುಟ್ಟು}}
\textbf{\kn{ಪಿಸುಗುಟ್ಟು}} — \emph{pisuguṭṭu} — \textquotedblleft{}to whisper\textquotedblright{}.

\begin{grammarlens}[title={What you've built}]
One more word for this chapter.
\end{grammarlens}

\subsection*{Your turn}
\begin{itemize}
\raggedright
  \item \emph{Say it:} \emph{pisuguṭṭu}
  \item \emph{Say it:} \emph{pisuguṭṭu}, once more
  \item \emph{Say it:} say \emph{pisuguṭṭu}
  \item \emph{Recall:} say the Kannada for to solve, then the Kannada for to dig, then say \emph{pisuguṭṭu} again
\end{itemize}

\begin{tcolorbox}[breakable,colback=teal!4,colframe=teal!35!black,title={Before you move on}]
What does \kn{ಪಿಸುಗುಟ್ಟು} mean? (\textquotedblleft{}To whisper\textquotedblright{}.) Say it once more.
\end{tcolorbox}
\section[\texorpdfstring{\kn{ಚಪ್ಪಾಳೆ} \kn{ತಟ್ಟು} (cappāḷe taṭṭu)}{ಚಪ್ಪಾಳೆ ತಟ್ಟು (cappāḷe taṭṭu)}]{\kn{ಚಪ್ಪಾಳೆ} \kn{ತಟ್ಟು} (cappāḷe taṭṭu) — to clap}

\label{lesson:KA-C280-cappale-tattu}



\begin{quote}
Before the new one: say the Kannada for to dig, then the Kannada for to whisper.
\end{quote}

\subsection*{You'll want to know: \kn{ಚಪ್ಪಾಳೆ} \kn{ತಟ್ಟು}}
\textbf{\kn{ಚಪ್ಪಾಳೆ} \kn{ತಟ್ಟು}} — \emph{cappāḷe taṭṭu} — \textquotedblleft{}to clap\textquotedblright{}.

\begin{grammarlens}[title={What you've built}]
One more word for this chapter.
\end{grammarlens}

\subsection*{Your turn}
\begin{itemize}
\raggedright
  \item \emph{Say it:} \emph{cappāḷe taṭṭu}
  \item \emph{Say it:} \emph{cappāḷe taṭṭu}, once more
  \item \emph{Say it:} say \emph{cappāḷe taṭṭu}
  \item \emph{Recall:} say the Kannada for to dig, then the Kannada for to whisper, then say \emph{cappāḷe taṭṭu} again
\end{itemize}

\begin{tcolorbox}[breakable,colback=teal!4,colframe=teal!35!black,title={Before you move on}]
What does \kn{ಚಪ್ಪಾಳೆ} \kn{ತಟ್ಟು} mean? (\textquotedblleft{}To clap\textquotedblright{}.) Say it once more.
\end{tcolorbox}
\section[\texorpdfstring{\kn{ಕೈ} \kn{ಕುಲುಕು} (kai kuluku)}{ಕೈ ಕುಲುಕು (kai kuluku)}]{\kn{ಕೈ} \kn{ಕುಲುಕು} (kai kuluku) — to shake hands}

\label{lesson:KA-C280-kai-kuluku}



\begin{quote}
Before the new one: say the Kannada for to whisper, then the Kannada for to clap.
\end{quote}

\subsection*{You'll want to know: \kn{ಕೈ} \kn{ಕುಲುಕು}}
\textbf{\kn{ಕೈ} \kn{ಕುಲುಕು}} — \emph{kai kuluku} — \textquotedblleft{}to shake hands\textquotedblright{}.

\begin{grammarlens}[title={What you've built}]
One more word for this chapter.
\end{grammarlens}

\subsection*{Your turn}
\begin{itemize}
\raggedright
  \item \emph{Say it:} \emph{kai kuluku}
  \item \emph{Say it:} \emph{kai kuluku}, once more
  \item \emph{Say it:} say \emph{kai kuluku}
  \item \emph{Recall:} say the Kannada for to whisper, then the Kannada for to clap, then say \emph{kai kuluku} again
\end{itemize}

\begin{tcolorbox}[breakable,colback=teal!4,colframe=teal!35!black,title={Before you move on}]
What does \kn{ಕೈ} \kn{ಕುಲುಕು} mean? (\textquotedblleft{}To shake hands\textquotedblright{}.) Say it once more.
\end{tcolorbox}
\section[\texorpdfstring{\kn{ತಲುಪಿಸು} (talupisu)}{ತಲುಪಿಸು (talupisu)}]{\kn{ತಲುಪಿಸು} (talupisu) — to deliver}

\label{lesson:KA-C280-talupisu}



\begin{quote}
Before the new one: say the Kannada for to clap, then the Kannada for to shake hands.
\end{quote}

\subsection*{You'll want to know: \kn{ತಲುಪಿಸು}}
\textbf{\kn{ತಲುಪಿಸು}} — \emph{talupisu} — \textquotedblleft{}to deliver\textquotedblright{}.

\begin{grammarlens}[title={What you've built}]
One more word for this chapter.
\end{grammarlens}

\subsection*{Your turn}
\begin{itemize}
\raggedright
  \item \emph{Say it:} \emph{talupisu}
  \item \emph{Say it:} \emph{talupisu}, once more
  \item \emph{Say it:} say \emph{talupisu}
  \item \emph{Recall:} say the Kannada for to clap, then the Kannada for to shake hands, then say \emph{talupisu} again
\end{itemize}

\begin{tcolorbox}[breakable,colback=teal!4,colframe=teal!35!black,title={Before you move on}]
What does \kn{ತಲುಪಿಸು} mean? (\textquotedblleft{}To deliver\textquotedblright{}.) Say it once more.
\end{tcolorbox}
\section[\texorpdfstring{\kn{ಹಿಂದಿಕ್ಕು} (hindikku)}{ಹಿಂದಿಕ್ಕು (hindikku)}]{\kn{ಹಿಂದಿಕ್ಕು} (hindikku) — to overtake}

\label{lesson:KA-C280-hindikku}



\begin{quote}
Before the new one: say the Kannada for to shake hands, then the Kannada for to deliver.
\end{quote}

\subsection*{You'll want to know: \kn{ಹಿಂದಿಕ್ಕು}}
\textbf{\kn{ಹಿಂದಿಕ್ಕು}} — \emph{hindikku} — \textquotedblleft{}to overtake\textquotedblright{}.

\begin{grammarlens}[title={What you've built}]
One more word for this chapter.
\end{grammarlens}

\subsection*{Your turn}
\begin{itemize}
\raggedright
  \item \emph{Say it:} \emph{hindikku}
  \item \emph{Say it:} \emph{hindikku}, once more
  \item \emph{Say it:} say \emph{hindikku}
  \item \emph{Recall:} say the Kannada for to shake hands, then the Kannada for to deliver, then say \emph{hindikku} again
\end{itemize}

\begin{tcolorbox}[breakable,colback=teal!4,colframe=teal!35!black,title={Before you move on}]
What does \kn{ಹಿಂದಿಕ್ಕು} mean? (\textquotedblleft{}To overtake\textquotedblright{}.) Say it once more.
\end{tcolorbox}
\section[(dialogue)]{Words from chapters 272-275 — a first pass}

\label{lesson:KA-R280-first-pass-words-from-chapters-272-275}



\begin{quote}
Say the words for to deliver and to overtake.
\end{quote}

\subsection*{Your turn}
Say these aloud:

\begin{itemize}
\raggedright
  \item to spill, to pour, to mix, to stir, to roast
  \item to fly, to wander around, to cross, to drive, to slip
  \item to dance, to press, to deposit, to withdraw, to sign
  \item to charge, to check, to translate, to apply, to compare
\end{itemize}

\begin{tcolorbox}[breakable,colback=teal!4,colframe=teal!35!black,title={Before you move on}]
To spill? (\textbf{\kn{ಚೆಲ್ಲು}}, \emph{cellu}.) To pour? (\textbf{\kn{ಸುರಿ}}, \emph{suri}.) To mix? (\textbf{\kn{ಬೆರೆಸು}}, \emph{beresu}.) To stir? (\textbf{\kn{ಕಲಕು}}, \emph{kalaku}.) To roast? (\textbf{\kn{ಹುರಿ}}, \emph{huri}.) To fly? (\textbf{\kn{ಹಾರು}}, \emph{hāru}.) To wander around? (\textbf{\kn{ಸುತ್ತಾಡು}}, \emph{suttāḍu}.) To cross? (\textbf{\kn{ದಾಟು}}, \emph{dāṭu}.) To drive? (\textbf{\kn{ಓಡಿಸು}}, \emph{ōḍisu}.) To slip? (\textbf{\kn{ಜಾರು}}, \emph{jāru}.) To dance? (\textbf{\kn{ಕುಣಿ}}, \emph{kuṇi}.) To press? (\textbf{\kn{ಒತ್ತು}}, \emph{ottu}.) To deposit? (\textbf{\kn{ಜಮಾ} \kn{ಮಾಡು}}, \emph{jamā māḍu}.) To withdraw? (\textbf{\kn{ಹಿಂಪಡೆ}}, \emph{himpaḍe}.) To sign? (\textbf{\kn{ಸಹಿ} \kn{ಮಾಡು}}, \emph{sahi māḍu}.) To charge? (\textbf{\kn{ಚಾರ್ಜ್} \kn{ಮಾಡು}}, \emph{cārj māḍu}.) To check? (\textbf{\kn{ಪರಿಶೀಲಿಸು}}, \emph{pariśīlisu}.) To translate? (\textbf{\kn{ಅನುವಾದಿಸು}}, \emph{anuvādisu}.) To apply? (\textbf{\kn{ಅರ್ಜಿ} \kn{ಹಾಕು}}, \emph{arji hāku}.) To compare? (\textbf{\kn{ಹೋಲಿಸು}}, \emph{hōlisu}.)
\end{tcolorbox}
\section[(dialogue)]{Words from chapters 276-280 — a second pass}

\label{lesson:KA-R280-second-pass-words-from-chapters-276-280}



\begin{quote}
Say the words for to scold and to tease.
\end{quote}

\subsection*{Your turn}
Say these aloud:

\begin{itemize}
\raggedright
  \item to scold, to tease, to quarrel, to console, to be proud
  \item to guess, to seem, to practise, to paint, to download
  \item to lend, to give back, to come back, to prevent, to prepare
  \item to prick, to spit, to yawn, to solve, to dig
  \item to whisper, to clap, to shake hands, to deliver, to overtake
\end{itemize}

\begin{tcolorbox}[breakable,colback=teal!4,colframe=teal!35!black,title={Before you move on}]
To scold? (\textbf{\kn{ಬೈ}}, \emph{bai}.) To tease? (\textbf{\kn{ರೇಗಿಸು}}, \emph{rēgisu}.) To quarrel? (\textbf{\kn{ಜಗಳವಾಡು}}, \emph{jagaḷavāḍu}.) To console? (\textbf{\kn{ಸಮಾಧಾನ} \kn{ಮಾಡು}}, \emph{samādhāna māḍu}.) To be proud? (\textbf{\kn{ಹೆಮ್ಮೆಪಡು}}, \emph{hemmepaḍu}.) To guess? (\textbf{\kn{ಊಹಿಸು}}, \emph{ūhisu}.) To seem? (\textbf{\kn{ಅನಿಸು}}, \emph{anisu}.) To practise? (\textbf{\kn{ಅಭ್ಯಾಸ} \kn{ಮಾಡು}}, \emph{abhyāsa māḍu}.) To paint? (\textbf{\kn{ಬಣ್ಣ} \kn{ಹಚ್ಚು}}, \emph{baṇṇa haccu}.) To download? (\textbf{\kn{ಡೌನ್ಲೋಡ್} \kn{ಮಾಡು}}, \emph{ḍaunlōḍ māḍu}.) To lend? (\textbf{\kn{ಸಾಲ} \kn{ಕೊಡು}}, \emph{sāla koḍu}.) To give back? (\textbf{\kn{ಹಿಂತಿರುಗಿಸು}}, \emph{hintirugisu}.) To come back? (\textbf{\kn{ಹಿಂತಿರುಗು}}, \emph{hintirugu}.) To prevent? (\textbf{\kn{ತಡೆ}}, \emph{taḍe}.) To prepare? (\textbf{\kn{ತಯಾರಿಸು}}, \emph{tayārisu}.) To prick? (\textbf{\kn{ಚುಚ್ಚು}}, \emph{cuccu}.) To spit? (\textbf{\kn{ಉಗುಳು}}, \emph{uguḷu}.) To yawn? (\textbf{\kn{ಆಕಳಿಸು}}, \emph{ākaḷisu}.) To solve? (\textbf{\kn{ಬಗೆಹರಿಸು}}, \emph{bageharisu}.) To dig? (\textbf{\kn{ಅಗೆ}}, \emph{age}.) To whisper? (\textbf{\kn{ಪಿಸುಗುಟ್ಟು}}, \emph{pisuguṭṭu}.) To clap? (\textbf{\kn{ಚಪ್ಪಾಳೆ} \kn{ತಟ್ಟು}}, \emph{cappāḷe taṭṭu}.) To shake hands? (\textbf{\kn{ಕೈ} \kn{ಕುಲುಕು}}, \emph{kai kuluku}.) To deliver? (\textbf{\kn{ತಲುಪಿಸು}}, \emph{talupisu}.) To overtake? (\textbf{\kn{ಹಿಂದಿಕ್ಕು}}, \emph{hindikku}.)
\end{tcolorbox}
